% 本文件由 03_代码/p4_tables.py 自动生成，请勿手工修改。
% 数据源：06_支撑材料/p4_results.json

\begin{table}[htbp]
  \centering
  \caption{两个交付文件在 2025 年 2 月 1 日--12 月 31 日（334 天）的实际运行结果}
  \label{tab:p4-branch}
  \footnotesize
  \begin{tabular}{lrr}
    \toprule
    指标 & result4-2（计划不得修改） & result4-3（6/12/18 可调整） \\
    \midrule
    计划量（$0{:}00$ 定下的量）/kWh & 21,800,593 & 20,514,619 \\
    最终生效常规量/kWh & 21,800,593 & 21,111,818 \\
    上调量/kWh & --- & 597,199 \\
    紧急购电量/kWh & 105,603 & 52,149 \\
    \midrule
    弃电量/kWh & 2,753,256 & 1,941,508 \\
    计划费/元 & 14,119,897 & 13,187,623 \\
    调整费/元 & --- & 714,424 \\
    紧急费/元 & 643,263 & 291,979 \\
    合计费用/元 & \textbf{14,763,160} & \textbf{14,194,026} \\
    期初储电量/kWh & 7,283 & 6,534 \\
    期末储电量/kWh & 6,102 & 5,941 \\
    \bottomrule
  \end{tabular}
  \par\vspace{2pt}
  \begin{minipage}{.92\textwidth}\footnotesize
    注：两个分支都在\textbf{同一条附件 4 实际价格路径}上、按各自的规则独立回放并独立标定，费用一栏不可相加。两者都按实际电价结算，包括初始计划费；主计费口径下不允许下调，故 result4-3 的下调量为零。期末储电量用于说明跨日状态差异——只比费用而不比期末储量是不完整的。
  \end{minipage}
\end{table}
